\section{Name Notes}
\label{sec:name-notes}

A Name Note is an Ironwood output whose memo records a name transition by encoding the type of name change, the name, the associated Unified Address when applicable, the expiration value, and a predecessor reference.

\subsection{Memo encoding}

The Name Note memo is a 512-byte field. A parser removes its maximal trailing sequence of zero bytes. The remaining bytes MUST be UTF-8. The colon byte (\texttt{0x3a}) separates fields. The first field MUST be the exact ASCII bytes \texttt{ZNS}. Extra fields, missing fields, and empty required fields are rejected.

The name field MUST contain between 1 and 63 bytes, each of which is ASCII \texttt{a}--\texttt{z}, \texttt{0}--\texttt{9}, or hyphen (\texttt{0x2d}). The first and last bytes MUST NOT be hyphen.

The remaining bytes encode one transition in one of the following formats:

\begin{center}
\texttt{ZNS:claim:<name>:<ua>:<expires\_at>:<prev\_rcm>}\\
\texttt{ZNS:update:<name>:<ua>:<expires\_at>:<prev\_rcm>}\\
\texttt{ZNS:release:<name>:<ua>:none:<prev\_rcm>}
\end{center}

The \texttt{<expires\_at>} field is either the canonical ASCII decimal encoding of a Unix timestamp in whole seconds or the exact ASCII bytes \texttt{none}. The value \texttt{none} indicates that the registration has no fixed expiration. A \texttt{release} MUST encode \texttt{none} as its \texttt{expires\_at} value and the Unified Address of the binding being released as its \texttt{ua} value.

% TODO: Mint still encodes/parses release Name Notes as ZNS:release:<name>::none:<prev_rcm> (empty UA), so u in σ is empty. zns-verify matches this section.

The \texttt{prev\_rcm} field is the 64-character lowercase hexadecimal encoding of the predecessor Name Note's \texttt{rcm} value. The first claim uses 64 zeroes.

A Name Note parser MUST reject a memo that does not match one of these encodings.

\subsection{Transition representation}

To derive the ZNS-specific note-commitment inputs, $\mathsf{rcm}_\sigma$ and $\psi_\sigma$, interpret the encoded transition as the tuple

$$\sigma = (\alpha, n, u, e, p)$$

where:
\begin{itemize}
  \item $\alpha$: The exact ASCII bytes \texttt{claim}, \texttt{update}, or \texttt{release}.
  \item $n$: The raw bytes of the name field.
  \item $u$: The raw bytes of the \texttt{ua} field.
  \item $e$: The raw bytes of the \texttt{expires\_at} field.
  \item $p$: The 32 raw bytes decoded from \texttt{prev\_rcm}.
\end{itemize}

\subsection{Commitment inputs}

ZNS derives both commitment inputs from the same transition representation $\sigma$ using the same hash construction. Within the hash construction, the only difference is the ASCII derivation tag, \texttt{rcm} or \texttt{psi}. The resulting digests are then reduced into different Pallas fields.

To prevent ambiguous concatenation, define $\mathsf{LP}(x)$ as the four-byte unsigned little-endian byte length of $x$, followed by $x$ itself. Define the protocol domain tag $T$ as the 12 ASCII bytes \texttt{ZcashName/v1}. Quotation marks are not part of $T$.

For derivation tag $t$, where $t$ is the exact ASCII byte string \texttt{rcm} or \texttt{psi}, define:

$$H_t(\sigma) = \mathsf{BLAKE2b{-}512}( \mathsf{LP}(T) \parallel \mathsf{LP}(t) \parallel \mathsf{LP}(\alpha) \parallel \mathsf{LP}(n) \parallel \mathsf{LP}(u) \parallel \mathsf{LP}(e) \parallel p ).$$

Because $e$ is included in $\sigma$, both ZNS-specific commitment inputs cryptographically bind the \texttt{expires\_at} value recorded in the Name Note.

% TODO: zns-resolver's verify call does not yet pass expires_at into H_t.

The final $p$ is exactly 32 raw bytes and has no length prefix. Each $H_t$ invocation uses unkeyed BLAKE2b-512 with a 64-byte output and no additional personalization.

The two commitment inputs are then:

$$\mathsf{rcm}_\sigma = \mathsf{ToScalar}(H_{\mathtt{rcm}}(\sigma)), \quad \psi_\sigma = \mathsf{ToBase}(H_{\mathtt{psi}}(\sigma)).$$

The value $\mathsf{rcm}_\sigma$ is the Pallas scalar-field element used as the Ironwood note-commitment randomness.

The value $\psi_\sigma$ is the Pallas base-field element supplied to the Ironwood note commitment construction.

$\mathsf{ToScalar}$, $\mathsf{ToBase}$, and the field types are those defined for Ironwood by the Zcash Protocol Specification \cite{zcash-protocol,zip224}. Neither $\mathsf{rcm}_\sigma$ nor $\psi_\sigma$ is a raw 64-byte BLAKE2b digest; each is the field element obtained by reducing the digest produced with its respective derivation tag. The canonical encoding of either field element is its 32-byte little-endian field representation. ZNS derives both field elements from $\sigma$, rather than from the $rseed$ derivation used by ordinary ZIP-212 receiving \cite{zip212,zcash-protocol}.

Changing the domain tag, a derivation tag, field order, length width, byte order, hash parameters, or reduction rule changes the protocol version and requires new test vectors.

\subsection{Commitment verification}

ZNS changes the derivation of $\psi$ and $\mathsf{rcm}$ but does not otherwise change the standard note commitment construction. A verifier derives $\psi_\sigma$ and $\mathsf{rcm}_\sigma$ from the encoded transition, combines them with the remaining note components using the standard construction \cite{zcash-protocol,zip224}, and reconstructs the note commitment.

The Name Note passes commitment verification only if the reconstructed value equals the \field{cmx} recorded in the on-chain action.

This equality establishes only that the transition encoded in the memo is cryptographically bound to that output. It does not establish Mint authorization, user authorization, or the current registry state.

\subsection{Commitment test vector}

For a \texttt{claim} transition with \texttt{name=alice}, \texttt{expires\_at=none}, and a 32-byte zero predecessor. Key material is derived from the all-zero 32-byte seed under unified account $\mathtt{m/32'/133'/1'}$, diversifier index 0, external scope, mainnet:

{\footnotesize
\begin{center}
\begin{tabularx}{\textwidth}{@{}lY@{}}
\toprule
Value & Canonical little-endian bytes rendered as hexadecimal \\
\midrule
\texttt{g\_d} & \texttt{de4338f2ab9fd8300a3a1c20dd690ce27026c6001c295d7c641a067ce809b11e} \\
\texttt{pk\_d} & \texttt{6df609f5710f3b5deecd4ee4b8f0173b44af6cf8918ac00269526031ba628996} \\
$\psi_\sigma$ & \texttt{9f8a61b860c737d4564f12c635d654b843bc7115d9dc6cf6f09e409c81b8d13e} \\
$\mathsf{rcm}_\sigma$ & \texttt{daa928be21d0ec13b5dbb0244699dbfeba546c71591d24d7824db78e4670c504} \\
\bottomrule
\end{tabularx}
\end{center}
}

\noindent\texttt{ua}: \texttt{u1897y9pzw3zk6n9twtzu2z5kpkzw3hms2c54fpyv8lnr79m73ta}\\
\texttt{zljkk3veaxrtwncp66lf45p3f274xy2amqckx0sraje4v835yw8q0q}

Set the value to zero and $\rho$ to 32 bytes of \texttt{0x33}. The resulting \texttt{cmx} is \texttt{cc320736a0c1df1e4ffcee2b64aa73a9e6d06bb218e155a6fef422e1ecb1f70c}. An implementation that produces another value is non-conforming.

The complete memo encoding of this transition is:

{\footnotesize
\begin{flushleft}
\texttt{ZNS:claim:alice:u1897y9pzw3zk6n9twtzu2z5kpkzw3hms2c54fpyv8lnr79m73ta}\\
\texttt{zljkk3veaxrtwncp66lf45p3f274xy2amqckx0sraje4v835yw8q0q:none:}\\
\texttt{0000000000000000000000000000000000000000000000000000000000000000}
\end{flushleft}
}
